% Part: model-theory
% Chapter: lindstrom
% Section: lindstrom-proof

\documentclass[../../../include/open-logic-section]{subfiles}

\begin{document}

\olfileid[fr]{mod}{lin}{prf}

\olsection{Le théorème de Lindström}

\begin{lem}
\ollabel{lem:lindstrom}
Soit $!E\in L(\Lang{L})$, où
$\Lang{L}$ est fini. Supposons
qu'il existe $n\in\Nat$ tel que,
pour toutes !!{structure}s
$\Struct{M}$ et~$\Struct{N}$,
$\Struct{M}\equiv_n\Struct{N}$
et $\Struct{M}\models_L !E$
entraînent
$\Struct{N}\models_L !E$.
Alors~$!E$ est équivalent à un
!!{sentence}~$!D$ du premier ordre :
$\Mod(L){!E}=\Mod(F){!D}$.
\end{lem}

\begin{proof}
Fixons un tel~$n$. La
\olref[bas][pis]{prop:qr-finite}
montre que, dans un langage
relationnel fini, les énoncés
du premier ordre de rang de
quantification au plus~$n$
ont seulement un nombre fini
de classes d'équivalence logique.
Pour toute !!{structure}~$\Struct{M}$,
choisissons des représentants de
ces classes et notons
$!D_{\Struct{M}}$ la conjonction
de ceux qui sont vrais dans
$\Struct{M}$. Ainsi,
$\Struct{N}\models !D_{\Struct{M}}$
si et seulement si
$\Struct{N}\equiv_n\Struct{M}$.

Parmi les types ainsi obtenus,
retenons ceux des structures qui
satisfont~$!E$, puis prenons la
disjonction finie correspondante :
\[
!D=\bigvee
  \Setabs{!D_{\Struct{M}}}{
    \Struct{M}\models_L !E},
\]
où l'on ne garde qu'un représentant
par classe; la disjonction vide
désigne une contradiction. Si
$\Struct{N}\models !D$, elle
est $n$-équivalente à une
structure satisfaisant~$!E$;
l'hypothèse donne donc
$\Struct{N}\models_L !E$.
Réciproquement, si
$\Struct{N}\models_L !E$, son
type est l'un des disjoints de
$!D$, et $\Struct{N}\models !D$.
Les deux classes de modèles
coïncident.\footnote{Corrections éditoriales : les conjonctions et disjonctions portent sur un nombre fini de représentants modulo l'équivalence logique, et non sur toutes les formules syntaxiques. Nous distinguons aussi la classe de modèles de $!E$ dans la logique abstraite de celle de $!D$ dans la logique du premier ordre.}
\end{proof}

\begin{thm}[Théorème de Lindström]
\ollabel{thm:lindstrom}
Si une logique abstraite
\emph{normale} $\tuple{L,\models_L}$
possède la propriété de compacité
et la propriété de Löwenheim--Skolem
descendante dénombrable, alors
$\tuple{L,\models_L}\leq
\tuple{F,\models}$. Comme la
normalité donne l'inégalité
inverse, les deux logiques
sont équivalentes.\footnote{Précision éditoriale : nous rétablissons la normalité, omise dans l'énoncé source mais utilisée dans sa preuve pour la dépendance finie, le renommage, la relativisation et la représentation du premier ordre.}
\end{thm}

\begin{proof}
Par la dépendance finie, il
suffit de traiter un énoncé
$!E\in L(\Lang{L})$ dans un
langage relationnel \emph{fini}.
Le \olref{lem:lindstrom} ramène
la preuve à l'existence d'un
rang~$n$ tel que deux structures
$n$-équivalentes s'accordent
sur~$!E$. Pour un langage fini
et relationnel, le
\olref[bas][pis]{thm:b-n-f}
exprime cette équivalence par
le va-et-vient de rang~$n$,
$I_n(\emptyset,\emptyset)$, en incluant les constantes dans toutes les clauses atomiques de base.

Supposons par l'absurde qu'aucun
rang ne suffise. Pour tout
$n\in\Nat$, choisissons des
structures $\Struct{M}_n$ et
$\Struct{N}_n$ telles que
$I_n(\emptyset,\emptyset)$,
$\Struct{M}_n\models_L !E$ et
$\Struct{N}_n\not\models_L !E$.
Nous allons coder simultanément
deux structures et la hauteur
disponible du va-et-vient dans
un langage auxiliaire au plus
dénombrable.\footnote{Adaptation de la preuve : nous explicitons le codage par une relation $H_k$ pour chaque longueur usuelle et par une théorie dénombrable de hauteurs. La source condense ce mécanisme en un énoncé et une relation sur des suites. Nous utilisons chaque paire de modèles séparément pour les tests finis ; la réunion proposée dans la source peut également être construite en choisissant des copies disjointes hors des valeurs communes des constantes, avec les mêmes relations sur celles-ci.}

Ce langage comporte deux copies
de~$\Lang{L}$, des !!{predicate}s
unaires $U,V,T$ pour les deux
domaines et les hauteurs, une
relation d'ordre~$<$ sur~$T$,
des constantes $0,t^*$, et,
pour chaque $k\geq0$, un
!!{predicate}~$H_k$ à $1+2k$
places. L'énoncé
$H_k(t,\mathbf a,\mathbf b)$
signifie que les $k$-uplets
$\mathbf a\in U^k$ et
$\mathbf b\in V^k$ disposent
encore d'une hauteur~$t$
de va-et-vient.\footnote{Précision du codage : la hauteur diminue, mais les deux domaines comparés restent fixes. Dans la construction source, le même indice désigne la paire de modèles et le rang ; son remplacement par $n^*-k$ ne garantit pas que les tuples restent dans la paire choisie à l'indice $n^*$. Le codage ci-dessous sépare ces rôles.} Pour $k=0$,
$H_0(t)$ concerne les deux
uplets vides. Le renommage et
la relativisation fournissent
des énoncés $!E_U,!E_V$
signifiant que les structures
des côtés~$U,V$ satisfont~$!E$.

Considérons la théorie~$\Theta$
qui contient $!E_U$ et
$\lnot !E_V$, ainsi que les
clauses suivantes, toutes
représentables dans la logique
normale par des énoncés du
premier ordre :
\begin{enumerate}
\item $U,V,T$ sont des parties non vides
  disjointes, avec $0,t^*\in T$; les constantes et
  relations des deux copies
  donnent des sous-structures
  sur~$U$ et~$V$.
\item $(T,<,0)$ est un ordre
  linéaire discret de premier
  élément~$0$, sans dernier
  élément; chaque hauteur possède
  un successeur immédiat et
  chaque hauteur non nulle un
  prédécesseur immédiat unique.
\item Pour chaque entier naturel
  usuel~$m$, $t^*$ se trouve
  au-delà du $m$-ième successeur
  de~$0$; cette condition est
  un schéma dénombrable d'%
  !!{sentence}s du premier ordre.
  On ajoute $H_0(t^*)$.
\item Pour tout~$k$ et toute
  hauteur~$t$, $H_k(t,\mathbf a,
  \mathbf b)$ impose $t\in T$, $\mathbf a\in U^k$ et $\mathbf b\in V^k$, ainsi que la
  correspondance des égalités,
  des constantes et des formules
  atomiques du langage fini.
  Si $t>0$ et~$s$ est son
  prédécesseur immédiat, il
  impose aussi les deux clauses
  de va-et-vient :
  \begin{align*}
  (\forall a\in U)(\exists b\in V)
    H_{k+1}(s,\mathbf a a,\mathbf b b),\\
  (\forall b\in V)(\exists a\in U)
    H_{k+1}(s,\mathbf a a,\mathbf b b).
  \end{align*}
\end{enumerate}
Chaque clause du dernier type
est du premier ordre. Comme
$\Lang{L}$ est fini, les
conditions atomiques pour une
longueur de tuple fixée sont
finies; $\Theta$ est donc
une théorie dénombrable.

Montrons qu'elle est finiment
satisfaisable. Une partie finie
ne demande que $t^*$ soit
au-delà d'un nombre fini de
successeurs de~$0$. Choisissons
$n$ plus grand que ces bornes,
prenons $T=\Nat$, $t^*=n$,
et plaçons les copies de
$\Struct{M}_n,\Struct{N}_n$
sur~$U,V$. Pour $t\leq n$,
interprétons $H_k(t,\mathbf a,
\mathbf b)$ par la relation
de va-et-vient de rang~$t$
sur ces $k$-uplets; pour
$t>n$, déclarons-la fausse.
Les définitions récursives
des relations de va-et-vient
vérifient toutes les clauses,
et $H_0(n)$ vaut par le choix
des deux structures. Cette
interprétation satisfait même
toutes les clauses de~$\Theta$
à l'exception éventuelle des
bornes $t^*>m$ non retenues.

Par compacité, $\Theta$ possède
un modèle~$\Struct{C}$. Sa
hauteur~$t^*$ a plus de
prédécesseurs que tout entier
usuel. Partant de
$t_0=t^*$, définissons, pour
chaque entier usuel~$k$,
$t_{k+1}$ comme le prédécesseur
immédiat de~$t_k$.
Aucun~$t_k$ n'est nul : sinon
$t^*$ serait le $k$-ième
successeur de~$0$, contrairement
aux bornes de~$\Theta$.
Pour chaque~$k$, retenons les
couples de $k$-uplets tels que
$H_k(t_k,\mathbf a,\mathbf b)$.
Ils décrivent des isomorphismes
partiels, contiennent la paire
vide par $H_0(t^*)$, et les
clauses de~$\Theta$ donnent
le va-et-vient de chaque paire
vers la longueur suivante.
Les deux sous-structures sur
$U,V$ sont donc partiellement
isomorphes. Par le
\olref[lsp]{thm:abstract-p-isom},
elles s'accordent sur~$!E$.
Mais $\Struct{C}\models_L !E_U$
et $\Struct{C}\models_L\lnot !E_V$,
contradiction. Le rang~$n$
recherché existe; le lemme
conclut la preuve.
\end{proof}

\end{document}
